\chapter{Data Understanding}
\label{ch:data-understanding}
\section{The source document}
Two candidate PDFs existed: a 66-page "MVP" export used in early exploration, and a 79-page "updated" export. The updated file was chosen because it alone contains the glossary, the abbreviations list, and the authority-code legend — reference data the agent needs for its own glossary and authority-code explanation features later on. This decision also expanded scope to include Chapter 3 (Non-Sovereign Operations), a full additional chapter of client relations, appraisal, trade finance, approval, and portfolio-management tables that the original \textasciitilde{}50-table MVP scope had not covered.

\section{Initial data exploration: what pdfplumber sees}
The project used \texttt{pdfplumber} for PDF text and character extraction. Early exploration established several facts about the document that shaped every downstream design decision:

\begin{itemize}
\item \textbf{Column headers are rotated 90 degrees.} Extracting them requires reading pdfplumber's raw characters in descending-top order, not the usual left-to-right reading order, and re-clustering headers that wrap across two vertical text strips (e.g. "Task Manager" printed as two adjacent runs) — verified character-by-character against real header text before any clustering logic was trusted.
\item \textbf{Authority codes carry superscript decorations that pdfplumber cannot always distinguish.} A cell reading "C13" (a Check/verify code with footnote 13) is rendered identically, in the underlying PDF, to a cell where "C" and "13" are two separate concepts — but pdfplumber's own word extraction is *inconsistent*: the same visual pattern comes out as one fused token on one page and two separate tokens on the next. This was confirmed at the character level (comparing raw x-position, y-position, and font size of each digit) before concluding that no text-only regex could solve it reliably — only real character geometry could, which the project subsequently built (\texttt{parsing/\allowbreak{}action\_\allowbreak{}geometry.py}).
\item \textbf{The letter "C" is semantically overloaded.} C1/C2 (printed in red) mean "Check and Verify"; C3/C4 (printed in purple) mean "Consult" (of the Board, Governors, or an independent function) — two unrelated authority concepts sharing one letter, distinguished only by level number and color. This was flagged early as a trap for the natural-language intent parser built later (Section 8.2), and is covered by a dedicated test ensuring a "consult" question can never surface Check/verify roles or vice versa.
\item \textbf{A hidden fourth hierarchy level exists.} Beyond chapter → process → task → child task, some child tasks carry their own lettered sub-items keyed by monetary threshold (e.g. activity 2.513.3 splits into "(a) Up to UA 2,000,000 / (b) Over UA 2M up to UA 10M / (c) Over UA 10M"). This was found via a single real screenshot and then confirmed document-wide (51 genuine instances, 2 false positives correctly excluded) before being formalized as a new node type, \texttt{threshold\_\allowbreak{}variant}.
\item \textbf{The informed-party marker "(i)" is not what it visually appears to be.} In the underlying character stream the marker is five separate characters — open-parenthesis, space, "i", space, close-parenthesis — not three consecutive characters as it looks on the rendered page. A first-generation extractor that scanned for three consecutive characters therefore missed every single instance; this was the single highest-impact data-understanding bug in the project (Section 7.4).
\item \textbf{Reference/redirect phrasing is inconsistent across the document.} At least four distinct real phrasings redirect a reader from one activity to another ("See DAM 16.100, 16.200...", "Refer to Activities X – Y in Section Z", "See 2.120: Organization of mission", and "(See 1.114.1 / 1.114.2)"), each discovered separately by hitting a real question that the agent answered wrongly, never assumed as a complete set in advance.
\end{itemize}
\section{Data quality issues catalogued before parsing began}
The project's discipline throughout was to confirm every anomaly against real page content (character coordinates, a screenshot, or both) before writing a fix — never to patch based on a hypothesis alone. This surfaced a long list of genuine data-quality issues in how pdfplumber represents the source PDF, distinct from any coding mistake:

\begin{itemize}
\item Superscript footnote digits are visually indistinguishable from level-indicator digits except by comparing each digit's own vertical alignment to the base letter's top/bottom edge — a rule only arrived at after an earlier, wrong hypothesis (comparing to cell-center) was tested against real coordinates and disproved.
\item A wide table row's rightmost and leftmost cells can render up to several points apart vertically, enough to be bucketed as two separate physical rows by naive y-coordinate grouping — confirmed via a requested screenshot of a real table (NSO section, page 58) after two more theoretical fixes had already been tried and rejected as unsafe.
\item Two adjacent but genuinely different columns (e.g. "Regional Implementation Support Manager (RISM)" and "Country Manager / DDG") can sit closer together in x-position than two wrapped lines of the *same* column elsewhere on the same page — meaning no single distance threshold can correctly separate "different columns" from "wrapped header lines" document-wide; a content-based signal (a closing parenthesis ending a balanced role name) had to be added alongside the geometric one.
\item pdfplumber's word extraction genuinely merges characters from two unrelated, visually overlapping lines of text into corrupted single "words" on a handful of titles — not synthetic OCR noise, but a real limitation of the library's row-heuristic on this specific document's layout, confirmed by re-extracting the same region at the character level and finding the real, un-merged text underneath.
\item One-off inconsistencies exist even within the source PDF itself: one footnote section header reads "Notes to" where every other instance in the document reads "Notes on" (a genuine typo in the DAM's own formatting), and it silently broke a section-boundary detector until checked page-by-page.
\end{itemize}
\section{Scope decisions grounded in real data}
Two scope-narrowing decisions were made only after checking the real content, not by default:

\begin{itemize}
\item The Abbreviations list (document pages 2–7) was extracted and wired into the agent, but the separate Glossary section (pages 8–11) was deliberately deferred: its longer, line-wrapping definitions break the simple row-pairing approach that works for the short Abbreviations entries, and it was judged not worth the additional geometry work within the time budget.
\item The action-geometry extractor's per-task action counts disagreed with the (already-trusted) text-based counts on roughly 14\% of tasks, with disagreements ranging from off-by-one to wildly wrong on individual cases. Rather than chase every remaining disagreement — the spread of error sizes suggested more than one uncatalogued cause — the project consciously shipped with the text-based count as the trusted source of truth and geometry as best-effort supplementary detail only, a decision made explicitly with the internship supervisor rather than silently absorbed into the code.
\end{itemize}
\section{Two deep-dive investigations worth narrating in full}
Two findings during Data Understanding are worth walking through in more procedural detail than the summary above, because the process by which they were found is as instructive as the findings themselves.

\textbf{The hidden fourth hierarchy level.} The DAM's hierarchy was originally understood as three levels — chapter, process, task — with an optional child task beneath a task. A screenshot of document page 36 (activity 2.513.3) showed something the schema had no place for: a child task with its own further sub-items, lettered by monetary threshold rather than numbered sequentially — "(a) Up to UA 2,000,000", "(b) Over UA 2M up to UA 10M", "(c) Over UA 10M" — each carrying its own independent set of authority-code responsibilities. The natural first question was whether this was a one-off formatting quirk on this single page or a real, recurring structural feature. It was checked document-wide before any schema change was made: 53 lettered sub-item lines were found across the full 79 pages, 51 genuine and 2 false positives (letter-lists embedded inside footnote prose, which the existing Notes-section boundary detector already excluded correctly). Confirmed as real and recurring, the decision — made jointly with the internship supervisor rather than unilaterally — was to model this as a new, distinct node type (\texttt{threshold\_\allowbreak{}variant}) with a dot-joined identifier built from whichever child task most recently opened (e.g. "2.513.3.a"), rather than treating it as just another child task, because it represents a conditional amount-band of one activity rather than a separate activity in its own right.

\textbf{The OCR-looking title corruption, root-caused rather than patched.} A small number of activity titles appeared to contain OCR-style letter-spacing artifacts (e.g. fragments resembling "ro", "je", "c" instead of readable words), and an early, reasonable-looking hypothesis was that this was scanning noise in the source PDF itself, to be handled with a cleanup pattern once convenient. Investigating properly instead of patching around the symptom revealed two distinct and unrelated root causes. The first, affecting one specific title, was a genuine parsing bug: the page's own footnote-section header read "Notes to NSO 3.510 – 3.520" where every other instance of that header in the entire document reads "Notes on" — a one-off typo in the DAM's own source formatting, confirmed by checking all 22 instances of the header document-wide (21 say "on", exactly 1 says "to") — and the section-boundary detector, which only recognized "on", consequently never closed that section and kept absorbing every following line, including the page footer, into the activity's title. The second, affecting the remaining cases, was not OCR noise at all: it was \texttt{pdfplumber}'s own word-extraction heuristic merging characters from two different, visually overlapping physical lines of real text into corrupted single tokens — confirmed by re-grouping the raw characters by row directly (bypassing \texttt{pdfplumber}'s own word-boundary logic entirely) and finding the real, uncorrupted two-line title underneath. The fix built a character-level row reconstruction used only as a targeted fallback when a title is flagged as corrupted, deliberately scoped narrowly rather than replacing the word-based pipeline everywhere, and reduced the known corruption count from six titles to one, with the residual case and the reasoning behind not chasing it further documented honestly rather than claimed as fully resolved.

